\documentclass[12pt,a4paper]{amsart}
\makeatletter
\renewcommand\normalsize{%
    \@setfontsize\normalsize{11.7}{14pt plus .3pt minus .3pt}%
    \abovedisplayskip 10\p@ \@plus4\p@ \@minus4\p@
    \abovedisplayshortskip 6\p@ \@plus2\p@
    \belowdisplayshortskip 6\p@ \@plus2\p@
    \belowdisplayskip \abovedisplayskip}
\renewcommand\small{%
    \@setfontsize\small{9.5}{12\p@ plus .2\p@ minus .2\p@}%
    \abovedisplayskip 8.5\p@ \@plus4\p@ \@minus1\p@
    \belowdisplayskip \abovedisplayskip
    \abovedisplayshortskip \abovedisplayskip
    \belowdisplayshortskip \abovedisplayskip}
\renewcommand\footnotesize{%
    \@setfontsize\footnotesize{8.5}{9.25\p@ plus .1pt minus .1pt}%%
    \abovedisplayskip 6\p@ \@plus4\p@ \@minus1\p@
    \belowdisplayskip \abovedisplayskip
    \abovedisplayshortskip \abovedisplayskip
    \belowdisplayshortskip \abovedisplayskip}
% ********************* DIMENSIONS:
% TEXT DIMENSIONS
\setlength\parindent    {30\p@}
\setlength\textwidth    {412\p@}
\setlength\textheight   {570\p@}
% PAPER AND TRIM SIZE
\paperwidth=210mm
\paperheight=260mm
\ifdefined\pdfpagewidth
\setlength{\pdfpagewidth}{\paperwidth}
\setlength{\pdfpageheight}{\paperheight}
\else
\setlength{\pagewidth}{\paperwidth}
\setlength{\pageheight}{\paperheight}
\fi
\calclayout
\makeatother






%\documentclass{amsart}

\usepackage{amsmath,amssymb}
\usepackage{mathrsfs}
\usepackage{mathtools}
\usepackage{stmaryrd}
\usepackage{enumerate}
\usepackage{tikz-cd}
\usepackage[all]{xy}
\usepackage{aliascnt}
\usepackage[colorlinks=true, linkcolor=blue, citecolor=blue]{hyperref}

\usepackage{enumitem}

\usepackage{combelow}

%\usepackage{fullpage}


\newtheorem{theorem}{Theorem}

\newaliascnt{lemma}{theorem}
\newtheorem{lemma}[lemma]{Lemma}
\aliascntresetthe{lemma}
\providecommand*{\lemmaautorefname}{Lemma}

\newaliascnt{corollary}{theorem}
\newtheorem{corollary}[corollary]{Corollary}
\aliascntresetthe{corollary}
\providecommand*{\corollaryautorefname}{Corollary}

\newaliascnt{proposition}{theorem}
\newtheorem{proposition}[proposition]{Proposition}
\aliascntresetthe{proposition}
\providecommand*{\propositionautorefname}{Proposition}

\newaliascnt{conjecture}{theorem}
\newtheorem{conjecture}[conjecture]{Conjecture}
\aliascntresetthe{conjecture}
\providecommand*{\conjectureautorefname}{Conjecture}

\newaliascnt{question}{theorem}
\newtheorem{question}[question]{Question}
\aliascntresetthe{question}
\providecommand*{\questionautorefname}{Question}



\theoremstyle{definition}

\newaliascnt{definition}{theorem}
\newtheorem{definition}[definition]{Definition}
\aliascntresetthe{definition}
\providecommand*{\definitionautorefname}{Definition}

\newaliascnt{remark}{theorem}
\newtheorem{remark}[remark]{Remark}
\aliascntresetthe{remark}
\providecommand*{\remarkautorefname}{Remark}

\newaliascnt{example}{theorem}
\newtheorem{example}[example]{Example}
\aliascntresetthe{example}
\providecommand*{\exampleautorefname}{Example}

\newaliascnt{notation}{theorem}
\newtheorem{notation}[notation]{Notation}
\aliascntresetthe{notation}
\providecommand*{\notationautorefname}{Notation}

\newtheorem*{conventions}{Conventions}
\newtheorem*{acknowledgements}{Acknowledgements}

%\newif\ifhascomments \hascommentstrue
%\ifhascomments
%  \newcommand{\matt}[1]{{\color{red}[[\ensuremath{\spadesuit\spadesuit\spadesuit} #1]]}}
%  \newcommand{\jeremy}[1]{{\color{red}[[\ensuremath{\clubsuit\clubsuit\clubsuit} #1]]}}
%\else
%  \newcommand{\matt}[1]{}
%  \newcommand{\jeremy}[1]{}
%\fi



\renewcommand{\setminus}{\smallsetminus}



\newcommand{\Gm}{\mathbb G_m}
\newcommand{\QT}{\mathbb Q}
\newcommand{\ZT}{\mathbb Z}
\newcommand{\RT}{\mathbb R}
\newcommand{\TT}{\mathbb T}


\newcommand{\Z}{\mathbb{Z}}
\newcommand{\ZZ}{\mathbb{Z}}
\newcommand{\bZ}{\mathbb{Z}}
\newcommand{\NN}{\mathbb{N}}
\newcommand{\N}{\mathbb{N}}
\newcommand{\QQ}{\mathbb{Q}}
\newcommand{\Q}{\mathbb{Q}}
\newcommand{\CC}{\mathbb{C}}
\newcommand{\R}{\mathbb{R}}
\newcommand{\RR}{\mathbb{R}}

\newcommand{\sep}{sep}

\newcommand{\cS}{\mathcal{S}}
\newcommand{\cT}{\mathscr{T}}
\newcommand{\cW}{\mathcal{W}}
\newcommand{\cX}{\mathcal{X}}
\newcommand{\cY}{\mathcal{Y}}
\newcommand{\cZ}{\mathcal{Z}}
%\newcommand{\sW}{\mathscr{W}\!}
\newcommand{\cI}{\mathcal{I}}
\newcommand{\cA}{\mathcal{A}}
\newcommand{\cC}{\mathcal{C}}
\newcommand{\cD}{\mathcal{D}}
\newcommand{\cU}{\mathcal{U}}
\newcommand{\cO}{\mathcal{O}}
\newcommand{\cV}{\mathcal{V}}
\newcommand{\cF}{\mathcal{F}}
\newcommand{\cG}{\mathcal{G}}
\newcommand{\cH}{\mathcal{H}}
%\newcommand{\cJ}{\mathcal{J}}
\newcommand{\cJ}{\mathscr{J}}
\newcommand{\mcJ}{\mathcal{J}}
\newcommand{\cK}{\mathcal{K}}
\newcommand{\cE}{\mathcal{E}}
\newcommand{\cP}{\mathcal{P}}
\newcommand{\cQ}{\mathcal{Q}}
\newcommand{\cL}{\mathcal{L}}
\newcommand{\cM}{\mathcal{M}}
\newcommand{\cR}{\mathcal{R}}
\newcommand{\cB}{\mathcal{B}}

\newcommand{\sL}{\mathscr{L}}
\newcommand{\sM}{\mathscr{M}}
\newcommand{\sI}{\mathscr{I}}
\newcommand{\sJ}{\mathscr{J}}
\newcommand{\sS}{\mathscr{S}}

\newcommand{\bL}{\mathbb{L}}
\newcommand{\bH}{\mathbb{H}}
\newcommand{\bG}{\mathbb{G}}
\newcommand{\bA}{\mathbb{A}}
\newcommand{\bP}{\mathbb{P}}

\newcommand{\fm}{\mathfrak{m}}
\newcommand{\g}{\mathfrak{g}}

\newcommand{\Var}{\mathbf{Var}}
\newcommand{\Stack}{\mathbf{Stack}}

\newcommand{\diff}{\mathrm{d}}
\newcommand{\red}{\mathrm{red}}
\newcommand{\id}{\mathrm{id}}
\newcommand{\Gor}{\mathrm{Gor}}

\newcommand{\str}{\mathrm{str}}


\renewcommand{\1}{\textbf{1}}
\newcommand{\<}{\langle}
\renewcommand{\>}{\rangle}
\DeclareMathOperator{\Lie}{Lie}
\DeclareMathOperator{\codim}{codim}
\DeclareMathOperator{\Vol}{Vol}
\DeclareMathOperator{\Stab}{Stab}
\DeclareMathOperator{\s}{s}
\DeclareMathOperator{\un}{un}
\DeclareMathOperator{\Bl}{Bl}
\DeclareMathOperator{\Sym}{Sym}
\DeclareMathOperator{\rep}{rep}
\DeclareMathOperator{\synt}{synt}
\DeclareMathOperator{\lci}{lci}
\DeclareMathOperator{\val}{val}
\DeclareMathOperator{\EP}{EP}
\DeclareMathOperator{\length}{len}
\DeclareMathOperator{\age}{age}
\DeclareMathOperator{\Jac}{Jac}
\DeclareMathOperator{\Res}{Res}
\DeclareMathOperator{\univ}{univ}
\DeclareMathOperator{\can}{can}
\DeclareMathOperator{\Art}{Art}
\DeclareMathOperator{\triv}{triv}
\DeclareMathOperator{\st}{st}
\DeclareMathOperator{\het}{ht}
\DeclareMathOperator{\nhet}{nht}
\DeclareMathOperator{\e}{e}
\DeclareMathOperator{\Spec}{Spec}
\DeclareMathOperator{\Ext}{Ext}
\DeclareMathOperator{\ord}{ord}
\DeclareMathOperator{\Mat}{Mat}
\DeclareMathOperator{\End}{End}
\DeclareMathOperator{\Hom}{Hom}
\DeclareMathOperator{\Isom}{Isom}
\DeclareMathOperator{\Aut}{Aut}
\DeclareMathOperator{\GL}{GL}
\DeclareMathOperator{\len}{len}
\DeclareMathOperator{\tors}{tors}
\DeclareMathOperator{\Tor}{Tor}
\DeclareMathOperator{\coh}{coh}
\DeclareMathOperator{\sm}{sm}
\DeclareMathOperator{\Fitt}{Fitt}
\DeclareMathOperator{\coker}{coker}
\DeclareMathOperator{\ordjac}{ordjac}
\DeclareMathOperator{\im}{im}
\DeclareMathOperator{\rank}{rank}
\DeclareMathOperator{\diag}{diag}
\DeclareMathOperator{\Div}{div}
\DeclareMathOperator{\SL}{SL}
\DeclareMathOperator{\Gr}{Gr}
\DeclareMathOperator{\uHom}{\underline{\Hom}}
\DeclareMathOperator{\uIsom}{\underline{\Isom}}
\DeclareMathOperator{\uAut}{\underline{\Aut}}
\DeclareMathOperator{\uSpec}{\underline{\Spec}}
\DeclareMathOperator{\Out}{Out}
\DeclareMathOperator{\Inn}{Inn}
\DeclareMathOperator{\Pic}{Pic}
\DeclareMathOperator{\HD}{HD}
\DeclareMathOperator{\Conj}{Conj}
\DeclareMathOperator{\inj}{inj}
\DeclareMathOperator{\rk}{rk}
\DeclareMathOperator{\Hdg}{Hdg}
\DeclareMathOperator{\omin}{omin}

\newcommand{\sT}{\mathscr{T}}
\newcommand{\sW}{\mathscr{W}}

\newcommand{\sHom}{\mathscr{H}om}

\DeclareMathOperator{\wt}{wt}



\newcounter{continuedCounter}







%For tikz-cd isomorphism arrow
\tikzset{cong/.style={draw=none,edge node={node [sloped, allow upside down, auto=false]{$\cong$}}},
         Isom/.style={above,every to/.append style={edge node={node [sloped, allow upside down, auto=false]{$\sim$}}}}}

%for resuming numbering in enumerate
\newcounter{ResumeEnumerate}










\title[Counter-examples to NCCR conjectures via Reductive McKay]{Counter-examples to Conjectures on NCCR\lowercase{s} via a Reductive McKay Correspondence}


\author{Matthew Satriano and Jeremy Usatine}

\thanks{MS was partially supported by an NSERC Grant}

\thanks{JU was partially supported by Simons Foundation MPS-TSM-00007918 and NSF DMS-2502347}

\address{Matthew Satriano, Department of Pure Mathematics, University of Waterloo}
\email{msatriano@uwaterloo.ca}

\address{Jeremy Usatine, Department of Mathematics, Florida State University}
\email{jusatine@fsu.edu}






























\begin{document}



\begin{abstract}
\vspace{-0.4em}
We disprove three conjectures appearing in Van den Bergh's plenary ICM survey, posed by De Deyn and Van den Bergh, which concern the interplay between non-commutative crepant resolutions (NCCRs) and stringy Euler characteristics. Our counter-examples are applications of a new McKay correspondence for reductive groups that we prove in this paper, extending the classical McKay correspondence for finite groups. We obtain explicit formulas for the stringy $E$-functions of reductive quotients in terms of the Weyl group action and stringy $E$-functions of toric varieties; we also obtain an equivalent formula expressed as an integral of the age function over Weyl-fixed loci of a compact real torus. For symmetric representations, our expressions are related to arithmetic Tutte polynomials, yielding combinatorial formulas for the stringy Euler characteristic.
\end{abstract}



\maketitle

\setcounter{tocdepth}{1}

%\tableofcontents



\numberwithin{theorem}{section}
\numberwithin{lemma}{section}
\numberwithin{corollary}{section}
\numberwithin{proposition}{section}
\numberwithin{conjecture}{section}
\numberwithin{question}{section}
\numberwithin{remark}{section}
\numberwithin{definition}{section}
\numberwithin{example}{section}
\numberwithin{notation}{section}


\vspace{-1.1em}

\section{Introduction}


The McKay correspondence began with a remarkable observation of John McKay relating the representation theory of finite groups to the minimal resolutions of Kleinian singularities. It has since blossomed into a far-reaching program at the intersection of algebraic geometry, representation theory, and mathematical physics, inspiring connections between the geometry of crepant resolutions, the representation theory of finite groups, and equivalences of derived categories.

Of particular interest in this paper is the non-commutative McKay correspondence, pioneered by Van den Bergh \cite{vdB2004} and inspired by the derived McKay correspondence of Bridgeland, King, and Reid \cite{BKR}. This rich theory connects non-commutative crepant resolutions (NCCRs) with geometric invariant theory, derived categories, maximal Cohen--Macaulay modules, tilting theory, and the minimal model program \cite{vdB10,IW13,SpenkoVanDenBergh17,vanDenBerghICM}.

The goal of this paper is twofold:
\begin{itemize}
\item Generalize the classical McKay correspondence from finite groups to reductive groups.
\item Apply our reductive McKay correspondence to disprove three conjectures concerning the relationship between NCCRs and stringy invariants.
\end{itemize}

Specifically, we disprove the following three conjectures, stated in Van den Bergh's plenary ICM survey \cite{vanDenBerghICM}. The first two were posed by De Deyn \cite{DeDeyn2023}.


\begin{conjecture}[{\cite[Conjecture 5.3]{vanDenBerghICM},\cite[Question 4.3]{DeDeyn2023}}]\label{conj:main1}
If $X$ is a normal Gorenstein variety over $\CC$ with an NCCR, then its stringy $E$-function is a polynomial.
\end{conjecture}

\begin{conjecture}[{\cite[Conjecture 5.12]{vanDenBerghICM},\cite[Question 4.11]{DeDeyn2023}}]\label{conj:main2}
The stringy Euler characteristic of a normal Gorenstein variety can be computed as the Euler characteristic of an NCCR, computed via periodic cyclic homology.
\end{conjecture}

\begin{conjecture}[{\cite[Conjecture 5.14]{vanDenBerghICM}}]\label{conj:main3}
Let $R$ be a normal Gorenstein ring which is either a complete local ring, or else connected $\NN$-graded (i.e., $R_0=\CC$). Assume $R$ has an NCCR $\End_R(M)$. Then the number of non-isomorphic indecomposable summands of $M$ is equal to $\chi_{\str}(X)$ for $X=\Spec R$.
\end{conjecture}

Our reductive McKay correspondences \autoref{thm:main-McKay-connected-reductive}, \autoref{thm:main-McKay-connected-reductive-2}, and \autoref{thm:main-McKay-connected-reductive-symmetric} give explicit representation-theoretic formulas for the stringy $E$-function and stringy Euler characteristic of reductive quotients. Our formulas are particularly amenable to computation and allow us to show the following theorem; see \autoref{thm:NCCR-counter-examples} for a more general statement.

\begin{theorem}
There exists a symmetric $\GL_2$-representation $V$ such that $V/\GL_2$ admits an NCCR and 
\[
	\chi_{\str}(V/\GL_2)\notin\ZZ
\]
In particular, \autoref{conj:main1}, \autoref{conj:main2}, and \autoref{conj:main3} are false.
\end{theorem}

To explain the origins of these conjectures and why they were expected to be true, we first recall the classical McKay correspondence and its connections with crepant resolutions as well as NCCRs.


\subsection{Background:~the McKay correspondence and NCCRs}
We begin by recalling McKay's original observation \cite{McKay}. Given a finite subgroup $\rho\colon G\hookrightarrow\SL_2(\CC)$, McKay considered the graph whose vertices are the irreducible representations $\rho_i$ of $G$ and the number of edges from $\rho_i$ to $\rho_j$ are given by the multiplicity of $\rho_i$ in the decomposition of $\rho_j\otimes\rho$. Surprisingly, as McKay showed, upon deleting the vertex corresponding to the trivial representation, this is precisely the dual graph of the exceptional curves in the minimal resolution of $\CC^2/G$.

Motivated by McKay's work as well as that of Gonz\'alez-Sprinberg and Verdier \cite{GonzalezSprinbergVerdier1983}, Reid \cite{Reid} proposed a far-reaching generalization of McKay's correspondence in higher dimensions. A resolution $\pi\colon X\to Y$ is called \emph{crepant} if $K_{X/Y}=0$. Reid conjectured that if $G\subset\SL_n(\CC)$ is a finite subgroup and $\pi\colon X\to\CC^n/G$ is a crepant resolution, then $\chi(X)$ is equal to the number of irreducible representations of $G$.

Reid's conjecture gave rise to an explosion of activity in the subject. Notably, it inspired Bridgeland, King, and Reid to show in their groundbreaking paper \cite{BKR} that the McKay Correspondence for threefolds lifts to an equivalence of derived categories. Shortly afterwards, Bondal and Orlov \cite{BO} conjectured that all crepant resolutions of a Gorenstein singularity are derived equivalent. This conjectural derived McKay correspondence led to a series of landmark results connecting derived categories to the minimal model program \cite{BrideglandFlop,BridgelandStability} and geometric invariant theory \cite{CI,HalpernLeistner,BallardFaveroKatzarkov}, as well as influencing developments in Gromov--Witten theory \cite{Ruan06,BG,CCIT,CCITtoric,CIJ} and Donaldson--Thomas theory \cite{BCY,DTCRC}.

The derived McKay correspondence was also the launching point for Van den Bergh's aforementioned non-commutative McKay correspondence. Van den Bergh \cite{vdB2004} proved that if 
$X$ has three-dimensional Gorenstein singularities and it has an NCCR $\Lambda$, then it also has a crepant resolution derived equivalent to $\Lambda$. This led to his broader conjecture that commutative and non-commutative crepant resolutions should be derived equivalent \cite{vdB2004}. As noted earlier, this was the starting point for deep interactions between non-commutative algebra, birational geometry, and representation theory \cite{vdB10,IW13,SpenkoVanDenBergh17,vanDenBerghICM}.

In parallel with the above geometric and categorical developments, a fundamentally different approach to the McKay correspondence arose through \emph{motivic integration}. The subject was pioneered by Kontsevich \cite{Kontsevich} and further developed by Denef and Loeser \cite{DenefLoeser1999,DenefLoeser2002}. Building on these ideas, Batyrev \cite{Batyrev1998} defined the \emph{stringy $E$-function} $E_{\str}(Y;u,v)$ which generalizes the Hodge--Deligne polynomial for smooth projective varieties and is invariant under crepant resolutions. Batyrev proved \cite{Batyrev99} that for a finite subgroup $G\subset\SL(V)$, we have
\begin{equation}\label{eqn:Batyrev-McKay}
	E_{\str}(V/G;u,v)=\sum_{g\in\Conj(G)}(uv)^{\dim V-\age_V(g)},
\end{equation}
where $\age_V(g)$ is a function determined by the eigenvalues of $g$; see \autoref{def:age}. In particular, if $X\to V/G$ is a crepant resolution, $E(X;u,v)=E_{\str}(V/G;u,v)$ and Batyrev's formula proved Reid's McKay correspondence.


\subsection{Stringy invariants and NCCRs:~the conjectures of De Deyn and Van den Bergh}
\autoref{conj:main1}, \autoref{conj:main2}, and \autoref{conj:main3} are incarnations of the philosophy that NCCRs should behave like usual crepant resolutions. Indeed, \autoref{conj:main1} is precisely the non-commutative analogue of the fact that if $X\to Y$ is a crepant resolution then $E_{\str}(Y)=E(X)$, and hence $E_{\str}(Y)$ is a polynomial. The Hochschild--Kostant--Rosenberg theorem \cite{HKR} essentially says that the cohomology of a smooth variety can be recovered from the periodic cyclic homology of its derived category. This suggests that the periodic cyclic homology of an NCCR should recover the stringy Euler characteristic \autoref{conj:main2}. Lastly, \autoref{conj:main3} is an incarnation of Reid's original McKay correspondence:~for a finite subgroup $G\subset\SL(V)$, Batyrev's formula \eqref{eqn:Batyrev-McKay} shows $\chi_{\str}(V/G)$ is the number of conjugacy classes of $G$, i.e., the number of irreducible representations of $G$, or equivalently, the number of indecomposable summands defining the canonical NCCR. \autoref{conj:main3} aims to generalize the phenomenon beyond the case of finite groups.

All of these conjectures are extremely natural and, in fact, there is substantial evidence in their favour. For a finite subgroup $G\subset\SL(V)$, \autoref{conj:main1} and \autoref{conj:main3} hold for the canonical NCCR of $V/G$ by Batyrev's formula \eqref{eqn:Batyrev-McKay} and \cite{Batyrev99}; see \cite[Examples 5.4 and 5.16]{vanDenBerghICM}. Batyrev \cite{Batyrev2000} additionally proved \autoref{conj:main1} for Gorenstein toric varieties. De Deyn proved \autoref{conj:main2} for the toric NCCRs constructed by \v{S}penko and Van den Bergh in \cite{SpenkoVanDenBergh2020i,SpenkoVanDenBergh2020ii}. Additionally, for the NCCRs constructed by Doyle in \cite{DeDeyn2023}, De Deyn proved \autoref{conj:main2} and \autoref{conj:main3}; see \cite[Example 5.17]{vanDenBerghICM}.

Another beautiful example comes from De Deyn's work \cite{DeDeyn2023} on $\Mat_{r\times n}/\SL_r$, where $\SL_r$ acts by left multiplication on $r\times n$ matrices. He showed that if $t=uv$, then
\[
	E_{\str}(\Mat_{r\times n}/\SL_r;u,v)={n\choose r}_{\! t}\frac{t^n(t-1)}{t^n-1},\qquad\chi_{\str}(\Mat_{r\times n}/\SL_r)=\frac{1}{n}{n\choose r},
\]
where ${n\choose r}_t$ is the $t$-binomial coefficient. In particular, the stringy $E$-function is a polynomial precisely when $\gcd(r,n)=1$. Remarkably, this is exactly the condition under which NCCRs were shown to exist by \v{S}penko--Van den Bergh \cite[Theorem 1.4.5]{SpenkoVanDenBergh17} and Doyle \cite[Theorem 1.1]{Doyle21}.





\subsection{Overview of reductive McKay and our counter-examples}
It has been a longstanding problem to obtain a representation-theoretic McKay correspondence outside the case of finite groups. Quotients by reductive groups occur naturally throughout algebraic geometry, particularly in geometric invariant theory and in the study of moduli spaces. In many of these settings, crepant resolutions are difficult to construct, while non-commutative crepant resolutions do exist thanks to the work of \cite{SpenkoVanDenBergh17}. This makes the lack of an analogous representation-theoretic McKay formula even more striking.

Recently, Groechenig, Wyss, and Ziegler \cite{GWZ} and Forey, Loeser, and Wyss \cite{FLW} made a major breakthrough in extending the McKay correspondence beyond the case of finite groups. Building on their earlier applications of $p$-adic integration to the fundamental lemma and mirror symmetry, Groechenig--Wyss--Ziegler proved an orbifold formula for quotients by reductive groups, expressing $p$-adic integrals in terms of cyclotomic inertia stacks. Forey--Loeser--Wyss then generalized these formulas to the motivic setting. 

In this paper, we build on their work to obtain an \emph{explicit representation-theoretic} McKay correspondence. We show that the stringy $E$-function of a reductive quotient $V/G$ can be computed directly from the Weyl group of $G$, its root system $\Phi$, and the weights of $V$, see \autoref{thm:main-McKay-connected-reductive} and \autoref{thm:main-McKay-connected-reductive-2}. Furthermore, in \autoref{thm:main-McKay-connected-reductive-symmetric}  we show a surprising connection with arithmetic Tutte polynomials \cite{BrandenMoci2014}.

Furthermore, our formulas also allow us to explain a phenomenon raised by Batyrev and Kiem. Noting that $E_{\str}(V/G)$ is a polynomial when $G\subset\SL(V)$ is a finite subgroup, Batyrev asked in \cite[Question 5.5]{Batyrev1998} whether this remains true for semi-simple $G$. In \cite{Kiem2003}, Kiem gave a counter-example. Our formulas in \autoref{thm:main-McKay-connected-reductive} explain the presence of denominators directly from the action of the Weyl group.

Given our explicit formula for $E_{\str}(V/G)$ and $\chi_{\str}(V/G)$, we had initially hoped to \emph{prove} \autoref{conj:main3} for quasi-symmetric representations. However, after working out several special instances of \autoref{thm:main-McKay-connected-reductive-2}, we found it was easy to produce non-polynomial stringy $E$-functions in this setting. This made us doubt \autoref{conj:main1}. After obtaining counter-examples to \autoref{conj:main1}, we saw they were counter-examples to \autoref{conj:main2} and \autoref{conj:main3} as well.\\
\\
We are now ready to state our main results.


\subsection{McKay Correspondence for linearly reductive groups}


Let $\Conj_{\inj}(\mu_m,G)$ denote the conjugacy classes of injective maps $\mu_m\to G$. Following \cite{Ekedahl}, for any finite type Artin stack $\cX$ over $\CC$, we let $\{\cX\}\in K_0(\Stack_\CC)$ be its associated class in the Grothendieck rings of stacks. Additionally for any finite type scheme $Y$ over $\CC$, we let $\{Y\}_{\str}$ denote the stringy motivic class of $Y$, see e.g., \cite[\S2.2]{SatrianoUsatine} and \cite[\S1.1]{SatrianoUsatine9} for background.

We recall the definition of the age function as well as the notion of a generic $G$-representation. The latter plays a key role in the work of \v{S}penko and Van den Bergh.

\begin{definition}\label{def:age}
Let $V$ be a $G$-representation and $\phi\colon\mu_m\to G$ an injective homomorphism. Choose integers $0\leq a_i<m$ such that the weights of the induced $\mu_m$-action on $V$ are given by $\overline{a}_i\in\ZZ/m$. Then the \emph{age} of $\phi$ is given by
\[
    \age_V(\phi):=\frac{1}{m}\sum_{i=1}^{\dim V} a_i.
\]
\end{definition}


\begin{definition}[{\cite[Definition 1.6]{SpenkoVanDenBergh17}}]
\label{def:generic-G-rep}
Given a finite-dimensional representation $V$ of a linear algebraic group $G$, let $U$ be the locus of points with closed orbit and trivial stabilizer. We say $V$ is a \emph{generic} $G$-representation if $\codim(V\setminus U)\geq2$.
\end{definition}

\begin{remark}
Let $V$ be a generic $G$-representation with $G\subset\SL(V)$ a linearly reductive subgroup. Then $V/G$ is $\QQ$-Gorenstein; if $G$ is connected or finite, then $V/G$ is Gorenstein; see \cite[Proposition 10.4]{SatrianoUsatine4}. If $G'\subset\SL_n(V')$ is a finite subgroup, then $V'$ is automatically a generic $G'$-representation.
\end{remark}



Our first result is a motivic McKay correspondence for quotients by linearly reductive groups; these groups need not be connected. This result is the basis for all of our later explicit representation-theoretic McKay correspondences.

\begin{theorem}\label{thm:main-lin-red-McKay}
Let $G$ be a linearly reductive group scheme and let $V$ be a generic $G$-representation with $G\subset\SL(V)$. Then
\[
	\{V/G\}_{\str}=-\bL^{\dim V-\dim G}\lim_{q\to\infty}\sum_{m\ge1}\ \sum_{\phi\in\Conj_{\inj}(\mu_m,G)} \bL^{-\age_V(\phi)+\age_{\g}(\phi)+\dim Z_G(\phi)}\{BZ_G(\phi)\}\sum_{n\ge1}q^{nm}.
\]
\end{theorem}


\begin{remark}
When $G$ is a finite group, \autoref{thm:main-lin-red-McKay} recovers the classical motivic McKay correspondence. Indeed, in this case the first two sums are finite, $\dim G=\dim Z_G(\phi)=\age_\g(\phi)=0$,
and the limit is $-1$. Thus, after identifying $\Conj_{\inj}(\mu_m,G)$ with conjugacy classes of elements of order $m$, the above formula becomes
\[
	\{V/G\}_{\str}=\sum_{g\in\Conj(G)}\bL^{\dim V-\age_V(g)}\{BZ_G(g)\}.
\]
\end{remark}


\subsection{Representation-theoretic McKay Correspondence for reductive groups}

For the remainder of the introduction, we fix the following notation for connected reductive $G$. Let $T\subset G$ be a choice of maximal torus, and let $W=N_G(T)/T$ be the Weyl group. Let $M=X^*(T)$ be the character lattice and $N$ be the cocharacter lattice of $T$. For every diagonalizable group scheme $H$, we let
\[
	\TT_H=\Hom(X^*(H),\RR/\ZZ).
\]
In particular, 
\[
	\TT:=\TT_T\simeq N_\RR/N
\]
which is a real torus. Note that $W$ acts on $T$, hence also on $\TT$; we let $\TT^w$ denote the locus of $\TT$ stabilized by $w\in W$. When $G$ is semi-simple and simply connected, $\TT$ agrees with the \emph{Steinberg torus} which is well-studied in Coxeter theory \cite{DilksPetersenStembridge2009,AguiarPetersen2015}.

Let $\beta_i$ denote the weights of $V$ listed with multiplicity. One may identify $\Conj_{\inj}(\mu_m,G)$ with the $W$-orbits of torsion points of $\TT$ of exact order $m$. Note that the fractional part $\{\beta_i\}$ gives a well-defined map $\TT\to[0,1)$, and if we define the function
\[
	\age_V\colon\TT\to\RR,\qquad \age_V(x):=\sum_i\{\beta_i(x)\}
\]
then:~(i) it is constant on $W$-orbits and, (ii) restricting this function to $W$-orbits of torsion points of $\TT$, it precisely recovers the usual definition of $\age_V$ on $\Conj_{\inj}(\mu_m,G)$. See \autoref{prop:TT-basics} for further details.

For any character $\lambda\in X^*(T)$, we let 
\[
	H_\lambda:=\{x\in\TT\mid \lambda(x)=0\in\RR/\ZZ\}.
\]
Letting $\Phi$ denote the roots of $G$ and $\Phi^+$ a choice of positive roots, we have a function $h_\g\colon\TT\to\NN$ given by
\[
	h_\g(x):=\#\{\alpha\in\Phi^+\mid x\in H_\alpha\}.
\]

With this notation, we obtain the following representation-theoretic version of the McKay correspondence for reductive groups. The integral in the statement of the theorem is an Euler-characteristic integral. Note the formal resemblance of this expression with an integral on the inertia stack of $[\TT/W]$.

\begin{theorem}[{McKay correspondence for connected reductive groups, I}]\label{thm:main-McKay-connected-reductive}
Let $G$ be a connected reductive group scheme and $V$ a generic $G$-representation with $G\subset\SL(V)$. If $t=uv$, then
\[
	E_{\str}(V/G;u,v)=\frac{t^{\dim V-|\Phi^+|}}{|W|}\sum_{w\in W}\frac{\int_{\TT^w}t^{-\age_V-h_\g} d\chi_c}{\det(t-w\mid M_\QQ)}.
\]
Furthermore, $\age_V$ takes only finitely many values; see \autoref{prop:TT-basics}(\ref{strata-cnst}).
\end{theorem}

\begin{remark}
As alluded to earlier, our formulas explain the presence of denominators, a phenomenon raised by Batyrev \cite[Question 5.5]{Batyrev1998} and Kiem \cite{Kiem2003}. The denominators arise precisely from the action of the Weyl group on the character lattice of a maximal torus.
\end{remark}




In fact, we may rewrite $E_{\str}(V/G)$ purely in terms of stringy $E$-functions of \emph{toric varieties}. To describe this, we introduce the following notation. For each $w\in W$, we let
\[
	\pi_w\colon M\to M_w:=M/(w-1)M,\qquad 	r_w:=\rank M_w.
\]
Note that $\rank M_w=\rank(M^w)$, and so $r_w$ is equal to the dimension of the $1$-eigenspace of $w$ acting on $M_\QQ$.

For any $S\subset M$, let
\[
	\rho_w(S):=\dim\QQ\pi_w(S),\qquad T_{w,S}=\Hom(M_w/\ZZ\pi_w(S),\bG_m)\subset T.
\]
Note that $T_{w,S}$ is a (possibly disconnected) diagonalizable group.

With this notation in place, we show $E_{\str}(V/G)$ may be written as an alternating weighted sum in terms of the $E_{\str}(V/T_{w,S})$.


\begin{theorem}[{McKay correspondence for connected reductive groups, II}]\label{thm:main-McKay-connected-reductive-2}
Let $G$ be a connected reductive group scheme and $V$ a generic $G$-representation with $G\subset\SL(V)$. Let $t=uv$ and for each $w\in W$, let
\[
	\Delta_w(t):=\det(t-w\mid M_\QQ/M_\QQ^w).
\]
Then
\[
	E_{\str}(V/G)=\frac{t^{-|\Phi^+|}}{|W|}\sum_{w\in W}\frac{1}{\Delta_w(t)}\sum_{R\subset\Phi^+} (-t^{-1})^{|R|}(t-1)^{|R|-\rho_w(R)} E_{\str}(V/T_{w,R}).
\]
Furthermore,
\[
	\chi_{\str}(V/G)=\frac{1}{|W|}\sum_{w\in W}\frac{1}{\Delta_w(1)}\sum_{\substack{R\subset\Phi^+\\|R|=\rho_w(R)}} (-1)^{|R|}\chi_{\str}(V/T_{w,R}).
\]
\end{theorem}

\begin{remark}
Note that the $w$-action on $M_Q / M^w_Q$ has trivial $1$-eigenspace, and so $Delta_w(1)$ does not vanish.
\end{remark}






\subsection{Symmetric representations via the arithmetic Tutte polynomial}
For symmetric representations, we obtain a surprising connection with the theory of arithmetic Tutte polynomials; the study of such polynomials is an active area at the intersection of combinatorics, matroids, and the theory of hyperplane arrangements \cite{Ardila07,ACH15,BrandenMoci2014,DupontFinkMoci,MaJinYang24}. Specifically, our integrals over $\TT^w$ can be understood as specializations of the \emph{multivariate arithmetic Tutte polynomial} $Z_\cA(q,\textbf{v})$, as introduced by Br\"and\'en and Moci in \cite{BrandenMoci2014}.

Let $\cW$ be the multi-set of weights of $V$ (counted with multiplicity). Since $V$ is symmetric, choosing a generic linear functional on $\ell\colon M_\RR\to\RR$, we can define $\cW^+\subset\cW$ to be the multi-subset of $\beta\in\cW$ with $\ell(\beta)>0$; note $\ell$ need not have any relation with the functional used to define $\Phi^+$. Let $\overline{\Phi}^+_w$ (resp.~$\overline{\cW}^+_w$) be the multi-set obtained by projecting $\Phi^+$ (resp.~$\cW^+$) via $\pi_w\colon M\to M_w$. (We use the bar so as not to conflict with standard Coxeter theory notation $\Phi^+_w$.) 

We now introduce the following specialization of multivariate arithmetic Tutte polynomial.

\begin{definition}\label{def:specialization-mat-poly}
For any $B\subset\cW^+$ and $R\subset\Phi^+$, let
\[
	m_w(B,R):=|( M_w/\ZZ\pi_w(B\cup R))_{\tors}|,\qquad \rho_w(B,R)=\dim\QQ\pi_w(B\cup R)
\]
and
\[
	\cM_w(q,x,y):=\sum_{\substack{B\subset\cW^+\\ R\subset\Phi^+}} m_w(B,R)q^{r_w-\rho_w(B,R)}x^{|B|}y^{|R|},
\]
which is a three-variable specialization of multivariate arithmetic Tutte polynomial.\footnote{Let $E_w$ be the arithmetic matroid given by the multi-set of vectors $\pi_w(\lambda)$ in $ M_w$ for all $\lambda$ in the mutli-set $\cW^+\sqcup\Phi^+$. Then $\cM_w(q,x,y)$ is obtained from $q^{r_w}Z_{E_w}(q,\textbf{v})$ by setting $v_e=x$ for $e\in\pi_w(\cW^+)$ and $v_e=y$ for $e\in\pi_w(\Phi^+)$.}
\end{definition}

With this definition in place, we prove:

\begin{proposition}\label{prop:main-intTw-arithmetic-Tutte}
Let $G$ be a connected reductive group scheme and $V$ a symmetric generic $G$-representation with $G\subset\SL(V)$. If $t=uv$, then
\[
	\int_{\TT^w}t^{-\age_V-h_\g}d\chi_c=t^{-|\cW^+|}\cM_w(0,t-1,t^{-1}-1).
\]
Furthermore, 
\[
	\widetilde{\cM}_w(t):=\frac{\cM_w(0,t-1,t^{-1}-1)}{(t-1)^{r_w}}\in\ZZ[t^\pm].
\]
\end{proposition}

Combining \autoref{thm:main-McKay-connected-reductive} and \autoref{prop:main-intTw-arithmetic-Tutte}, we have:


\begin{theorem}[{McKay correspondence for symmetric representations}]\label{thm:main-McKay-connected-reductive-symmetric}
Let $G$ be a connected reductive group scheme and $V$ a symmetric generic $G$-representation with $G\subset\SL(V)$. Let $t=uv$ and for each $w\in W$, let
\[
	\Delta_w(t):=\det(t-w\mid M_\QQ/M_\QQ^w).
\]
Then
\[
	E_{\str}(V/G;u,v)=\frac{t^{\dim V-|\Phi^+|-|\cW^+|}}{|W|}\sum_{w\in W}\frac{\widetilde{\cM}_w(t)}{\Delta_w(t)}.
\]
In particular, the stringy Euler characteristic is given by
\[
	\chi_{\str}(V/G)=\frac{1}{|W|}\sum_{w\in W}\frac{\widetilde{\cM}_w(1)}{\Delta_w(1)}.
\]
Explicitly, 
\[
	\widetilde{\cM}_w(1)=\!\!\!\!\!\!\!\!\!\sum_{\substack{B\subset\cW^+,\, R\subset\Phi^+\\ |B|+|R|=\rho_w(B,R)=r_w}}\!\!\!\!\!\!\!\!\!\!\!\!(-1)^{|R|}m_w(B,R).
\]
\end{theorem}



\subsection{Purity of crepant resolutions of group quotients}

As a first application of our McKay correspondences, we generalize Batyrev's purity theorem \cite[Theorem 8.4]{Batyrev99} from the case of finite group quotients to linearly reductive quotients.



\begin{theorem}\label{thm:purity-Hodge-Tate}
Let $G$ be a linearly reductive group over $\CC$ and let $V$ be a generic $G$-representation with $G\subset\SL(V)$. If $\pi\colon X\to V/G$ is a projective crepant resolution, then $H^i(X,\QQ)$ is pure of weight $i$ and is Hodge--Tate, i.e., for $d_i = \dim H^{2i}(X,\QQ)$, we have
\[
 H^{2i+1}(X,\QQ)=0,\qquad H^{2i}(X,\QQ)\simeq\QQ(-i)^{\oplus d_i}.
\]
\end{theorem}





\subsection{Counter-examples to conjectures of De Deyn and Van den Bergh for NCCRs}

Lastly, we apply our McKay correspondences to disprove \autoref{conj:main1}, \autoref{conj:main2}, and \autoref{conj:main3}.

%Note that all three of these conjectures imply that if an NCCR of $V/G$ exists, then $\chi_{\str}(V/G)\in\ZZ$. Making use of \autoref{thm:main-McKay-connected-reductive-symmetric}, we obtain the following theorem which gives a counter-example to all of the above conjectures.


\begin{theorem}\label{thm:NCCR-counter-examples}
Let $V_{(a,b)}$ denote the irreducible $\GL_2$-representation with highest weight vector $(a,b)$. Let $U:=V_{(1,0)}\oplus V_{(1,0)}^*$, $m\geq3$, and $a>b\geq0$. Then
\[
	V:=U^{\oplus m}\oplus (V_{(a,b)}\oplus V_{(a,b)}^*)
\]
is a generic symmetric representation. Furthermore, 
\[
	\chi_{\str}(V/\GL_2)\notin\ZZ\ \Longleftrightarrow\ a+b\equiv2\pmod{4}.
\]
When $(m,a,b)=(4,2,0)$, $V/\GL_2$ has an NCCR, hence it is a counter-example to \autoref{conj:main1}, \autoref{conj:main2}, and \autoref{conj:main3}.
\end{theorem}






\vspace{1em}


\noindent\textbf{Acknowledgments.} It is a pleasure to thank Wentang Kuo, Oliver Pechenik, David Speyer, and Hunter Spink. We especially wish to thank Vic Reiner and Ben Webster for their insights and for answering numerous technical questions on Coxeter theory. We are very grateful to Jim Geelen for suggesting a connection with the Tutte polynomial.

\vspace{1em}

\noindent\emph{AI disclosure.} The counter-examples were discovered by us. All ideas and arguments are due to the authors with one exception:~after Jim Geelen suggested a connection with Tutte polynomials, AI provided a proof of the special case of \autoref{prop:main-intTw-arithmetic-Tutte} when we formally set $\age_V=0$.














\section{McKay correspondence for linearly reductive groups}


In this section, we prove \autoref{thm:main-lin-red-McKay}. 



\begin{proposition}\label{prop:cyc-inertia-VmodG}
Let $G$ be a smooth affine algebraic group acting on a scheme $X$ over an algebraically closed field $k$ of characteristic $0$. Then for all $N\geq1$, we have
\[
	\uHom(B\mu_N,[X/G]) \simeq \coprod_{\phi\in\Conj(\mu_N,G)} [X^\phi/Z_G(\phi)].
\]
\end{proposition}
\begin{proof}
The same exact proof as in \cite[Proposition 10.2]{HuangSatrianoUsatine} shows that the $S$-valued points of $\uHom(B\mu_N,[X/G])$ are equivalent to the groupoid of triples $(P,\tau,\lambda)$ with $P\to S$ a $G$-torsor, $\tau\colon P\to X$ a $G$-equivariant map, $\lambda\colon\mu_N\times P\to P$ defining a group action, and where conditions (1), (2), (3), and (5) hold; note that condition (4) precisely corresponds to representability of the associated map $B\mu_{N,S}\to[X/G]$.

The remainder of the proof is nearly identical to that of \cite[Theorem 10.1]{HuangSatrianoUsatine}. Given $\phi\colon\mu_N\to G$, we construct a map
\[
    \alpha_\phi\colon [X^\phi/Z_G(\phi)]\to \uHom(B\mu_N,[X/G])
\]
as follows. An $S$-valued point of $[X^\phi/Z_G(\phi)]$ is given by $(Q,f)$ where $Q\to S$ is a left $Z_G(\phi)$-torsor and $f\colon Q\to X^\phi$ is a $Z_G(\phi)$-equivariant map. Let $P=G\times^{Z_G(\phi)}Q\to S$ be the induced $G$-torsor. We let $\tau\colon P\to X$ and the $\mu_N$-action on $P$ be given by
\[
	\tau([g,q])=gf(q),\qquad\zeta\cdot[g,q]=[g\phi(\zeta)^{-1},q].
\]
By construction, conditions (1), (2), (3), and (5) hold.

Lemma 10.4 of \cite{HuangSatrianoUsatine} shows that $\Conj(\mu_N,G)$ is a finite set. We show that the map
\[
    \alpha\colon\coprod_{\phi\in\Conj(\mu_N,G)}[X^\phi/Z_G(\phi)]\to \uHom(B\mu_N,[X/G])
\]
induced by the $\alpha_\phi$ is an equivalence of stacks.

To see $\alpha$ is faithful, let $(P,f)\in[X^\phi/Z_G(\phi)](S)$; since we can look fppf locally, we can assume $P\simeq Z_G(\phi)_S$ is the trivial torsor. Let $x=f(1)\in X^\phi(S)$. Then $\alpha_\phi$ maps $(P,f)$ to the triple $(G_S,\tau_x,\lambda_\phi)$, where
\[
	\tau_x(g)=gx,\qquad\zeta\cdot g=g\phi(\zeta)^{-1}.
\]
An automorphism of $(P,f)$ is given by right multiplication by an element $h\in Z_G(\phi)(S)$; this is mapped by $\alpha_\phi$ to right multiplication by $h$ on $G_S$, which we may assume to be the identity map. Thus, $h=1$, which proves faithfulness.

To show $\alpha$ is full, let $(P,f)\in[X^\phi/Z_G(\phi)](S)$ and $(P',f')\in[X^{\phi'}/Z_G(\phi')](S)$ and let $\xi\colon\alpha_\phi(P,f)\xrightarrow{\simeq}\alpha_{\phi'}(P',f')$ be an isomorphism. Since $\alpha$ is faithful, we can look locally on $S$ where we can again assume $P$ and $P'$ are trivial. Then we have $x=f(1)\in X^\phi(S)$ and $x'=f'(1)\in X^{\phi'}(S)$. Since $\xi$ is given by right multiplication by some $g\in G(S)$ and $\xi$ commutes with the $\mu_N$-actions, we see
\[
	\phi(\zeta)^{-1}g=g\phi'(\zeta)^{-1}
\]
for all scheme-theoretic points $\zeta$ of $\mu_N$, i.e., $\phi'=g^{-1}\phi g$; hence, $\phi$ and $\phi'$ are conjugate, so we can assume $\phi=\phi'$. In particular, $g\in Z_G(\phi)(S)$. Since $x=gx'$, multiplication by $g$ yields an isomorphism $(P,f)\xrightarrow{\simeq}(P',f')$ which maps to $\xi$ under $\alpha$.

Lastly, to prove essential surjectivity, suppose $(P,\tau,\lambda)$ is an $S$-valued point of $\uHom(B\mu_N,[X/G])$. Since $\alpha$ is fully faithful, we can look locally, where we may assume $P\simeq G_S$ is the trivial torsor. Since the $\mu_N$-action commutes with the $G$-action, there is a homomorphism $\widetilde{\phi}\colon\mu_{N,S}\to G_S$ such that
\[
    \zeta\cdot g=g\widetilde{\phi}(\zeta)^{-1}
\]
By \cite[Lemma 10.4]{HuangSatrianoUsatine}, after possibly replacing $S$ by a further fppf cover and changing our trivialization, there is a homomorphism $\phi\colon\mu_N\to G$ such that $\widetilde{\phi}=\phi_S$. Letting $x=\tau(1)\in X(S)$ and using that $\tau$ is both $G$-equivariant and $\mu_N$-invariant, we have
\[
    x=\tau(1)=\tau(\zeta\cdot1)=\tau(\phi(\zeta)^{-1})=\phi(\zeta)^{-1}x.
\]
So, $x\in X^\phi(S)$ and so defines an $S$-valued point of $[X^\phi/Z_G(\phi)]$ corresponding to the trivial torsor and the map $f\colon Z_G(\phi)_S\to X^\phi$ with $f(1)=x$. We see then that $\alpha_\phi(Z_G(\phi)_S,f)\simeq (P,\tau,\lambda)$.
\end{proof}







\begin{proof}[Proof of \autoref{thm:main-lin-red-McKay}]
Let $\cX=[V/G]$, $Y=V/G$, and let $\pi\colon \cX\to V/G=:Y$ be the good moduli space map. We begin by showing that
\[
	\{V/G\}_{\str} =-\bL^{\dim V-\dim G}\lim_{q\to\infty} \sum_{N\geq1}[I_{\mu_N}\cX]^w q^N,
\]
where $[I_{\mu_N}\cX]^w$ is the weighted motivic class defined just before Corollary 5.1.5 in \cite{FLW}; note that (loc.~cit) uses the non-standard notation $I_{\mu_N}\cX$ to denote $\uHom(B\mu_N,\cX)$. The formula above follows from Theorem 5.1.4 of (loc.~cit) upon showing $\dim Y=\dim V-\dim G$, $|\omega_{Y^\natural}|=|\omega_{\can}|$, and Assumptions 5.1.1 of (loc.~cit) hold for $\pi$ when $D=0$. Condition (2) holds by Proposition 10.3 and the proof of Proposition 10.4 in \cite{SatrianoUsatine4}, which shows there is a dense open subset $U\subset Y$ such that $\pi^{-1}(U)\to U$ is an isomorphism, $\codim(Y\setminus U)\geq2$, and $\codim(\cX\setminus\pi^{-1}(U))\geq2$. Furthermore, by \cite[Theorem 1.21]{SatrianoUsatine4}, we see $Y$ is log-terminal. Conditions (1) and (3) then follow from \cite[Remark 5.1.2]{FLW}. Condition (4) follows from \cite[Proposition 10.2(2)]{SatrianoUsatine4}. Lastly, since $Y$ is log-terminal (hence $\QQ$-Gorenstein), and by the existence of our $U$ above, \cite[Lemma 6.1.2]{FLW} shows that $|\omega_{Y^\natural}|$ agrees with $|\omega_{\can}|$. The existence of our $U$ also shows $\dim Y=\dim V-\dim G$. 

Next, by \autoref{prop:cyc-inertia-VmodG},
\begin{align*}
	[I_{\mu_N}\cX]^w &= \sum_{\phi\in\Conj(\mu_N,G)} \bL^{-w(\phi)}\{[V^\phi/Z_G(\phi)]\}\\
	&= \sum_{\phi\in\Conj(\mu_N,G)} \bL^{\dim V^\phi-w(\phi)}\{BZ_G(\phi)\},
\end{align*}
where the second equality holds because $[V^\phi/Z_G(\phi)]\to BZ_G(\phi)$ is the total space of a vector bundle whose rank is $\dim V^\phi$.

By the definition of the function $w$ given in \cite[\S2.9.7]{FLW}, we see that for any $\mu_N$-representation $E$, we have $w(E)=\dim E^\phi+\age_E(\phi)$. Furthermore, by the description of the weight function for the tangent complex of a global quotient in \S2.9.7 of (loc.~cit), we see
 \begin{align*}
 w(\phi) &=w(V)-w(\g)= \dim V^\phi+\age_V(\phi)-\dim\g^\phi-\age_\g(\phi)\\
 &=\dim V^\phi+\age_V(\phi)-\dim Z_G(\phi)-\age_\g(\phi).
 \end{align*}

Lastly, every homomorphism $\mu_N\to G$ factors uniquely through its image $\mu_m\subset G$ for some $m$ dividing $N$. Thus, 
\begin{align*}
	\sum_{N\geq1}[I_{\mu_N}\cX]^wq^N &=\sum_{N\geq1}\sum_{\phi\in\Conj(\mu_N,G)} \bL^{-\age_V(\phi)+\age_\g(\phi)+\dim Z_G(\phi)}\{BZ_G(\phi)\}q^N\\
	&=\sum_{m\geq1}\sum_{\phi\in\Conj_{\inj}(\mu_m,G)} \bL^{-\age_V(\phi)+\age_\g(\phi)+\dim Z_G(\phi)}\{BZ_G(\phi)\}\sum_{n\geq1}q^{nm},
\end{align*}
which finishes the proof.
\end{proof}





\section{McKay correspondence for connected reductive groups}





\begin{proposition}\label{prop:molien}
	Let $H$ be a linearly reductive group with identity component $H^\circ$. If $T\subset H^\circ$ is a maximal torus with cocharacter lattice $N$ and $t=uv$, then $W_H=N_H(T)/T$ is finite and 
	\[
		 E(BH;u,v)=\frac{t^{\dim T-\dim H}}{|W_H|}\sum_{w\in W_H}\frac{1}{\det(t-w\mid  M_\QQ)}.
	\]
\end{proposition}
\begin{proof}
	Letting $W_{H^\circ}$ be the Weyl group of $H^\circ$, \cite[6.1.6.1]{DeligneHodgeIII} shows
	\[
		H^*(BH^\circ,\QQ)\simeq H^*(BT,\QQ)^{W_{H^\circ}}.
	\]
	By \cite[Lemma 2.4]{Brion2015}, we have a short exact sequence
	\[
		1\to W_{H^\circ}\to W_H\to \pi_0(H)\to 1,
	\]
	so, in particular, $W_H$ is finite. Since $H^*(BH,\QQ)=H^*(BH^\circ,\QQ)^{\pi_0(H)}$ (see e.g., \cite[9.1.1]{DeligneHodgeIII}) we therefore have
	\begin{equation}\label{BH-equals-W-invariants-of-BT}
		H^*(BH,\QQ)\simeq H^*(BT,\QQ)^{W_H}=(\Sym M_\QQ)^{W_H}.
	\end{equation}
	By \cite[9.1.1]{DeligneHodgeIII}, $H^*(BH,\QQ)$ vanishes in odd degrees and $H^{2d}(BH,\QQ)$ is pure of type $(d,d)$ for all $d$. So, by Poincar\'e duality \cite[Paragraph after Definition 3.11]{Tubach2025} we have
	\[
		E(BH;u,v)=t^{-\dim H}\sum_{d\geq0} h^{2d}(BH,\QQ)t^{-d}.
	\]
	Then by \eqref{BH-equals-W-invariants-of-BT} and Molien's formula,
	\begin{align*}
		E(BH;u,v) &=\frac{t^{-\dim H}}{|W_H|}\sum_{w\in W_H}\det(1-t^{-1}w\mid M_\QQ)^{-1}\\
		&=\frac{t^{\dim T-\dim H}}{|W_H|}\sum_{w\in W_H}\det(t-w\mid M_\QQ)^{-1}.\qedhere
	\end{align*}
\end{proof}



Recall our notation from the introduction:

\begin{definition}\label{def:hyperplanes-h-fnc}
	For any diagonalizable group scheme $H$, we let
	\[
		\TT_H:=\Hom(X^*(H),\RR/\ZZ).
	\]
	If $G$ is reductive with maximal torus $T$ and $\lambda\in X^*(T)$, we let
	\[
		H_\lambda := \{x\in\TT\mid \lambda(x)=0\in\RR/\ZZ\}.
	\]
	For any set of characters $S:=\{\lambda_1,\dots,\lambda_s\}$ and subset $I\subset S$, we let
	\[
		H_{S,I}:=\bigcap_{i\in I}H_{\lambda_i}\setminus\bigcup_{j\notin I}H_{\lambda_j}
	\]
	and refer to the connected components of the $H_{S,I}$ as the strata determined by $S$.
\end{definition}


\begin{proposition}\label{prop:TT-basics}
Let $V$ be a representation of a reductive group $G$. Let $T\subset G$ be a maximal torus whose cocharacter lattice is $N$ and let $\TT:=\TT_T$. Let $\Phi^+$ be the positive roots of $G$.
\begin{enumerate}
	\item\label{can-bij} We have a canonical $W$-equivariant map
	\[
		\TT[m]\simeq\Hom(\mu_m,T)\to \Hom(\mu_m,G)
	\]
	which induces a canonical bijection
	\[
		\TT[m]/W\xrightarrow{\sim}\Conj(\mu_m,G).
	\]
	Under this bijection, the $W$-orbits of points with exact order $m$ map bijectively with $\Conj_{\inj}(\mu_m,G)$. Given $x\in\TT_{\tors}$, there is therefore a unique corresponding $\phi_x\in\Conj_{\inj}(\mu_m,G)$, where $m$ is the order of $x$.
	\item\label{WxNZ} If $x\in\TT_{\tors}$, then $T$ is a maximal torus of the identity component of $Z(\phi_x)$ and
	\[
		W_x:=\Stab_W(x)=N_{Z(\phi_x)}(T)/T.
	\]
	\item\label{age-extension} If $\beta_1,\dots,\beta_n$ are the weights of $V$ listed with multiplicity, and if we define
	\[
		\age_V\colon\TT\to\RR,\qquad \age_V(x):=\sum_i\{\beta_i(x)\},
	\]
	then $\age_V$ is constant on $W$-orbits, and $\age_V(x)=\age_V(\phi_x)$ for $x\in\TT_{\tors}$.
	\item\label{strata-cnst} If $G\subset\SL(V)$, and $S$ denotes the union of $\Phi^+$ and the set of weights of $V$, then $\age_V$ and $\age_\g$ are integer-valued and constant on the strata determined by $S$.
	\item\label{hVvsage} Let $V$ be a symmetric representation whose non-zero weights are $\pm\beta_1,\dots,\pm\beta_s$ listed with multiplicity and let
	\[
		h_V\colon\TT\to\NN,\qquad h_V(x):=\#\{i\mid \beta_i(x)=0\in\RR/\ZZ\}.
	\]
	Then for $x\in\TT$, we have
	\[
		\age_V(\phi_x)=s-h_V(x).
	\]
	In particular,
	\[
		\age_\g(\phi_x)=|\Phi^+|-h_\g(x).
	\]
\end{enumerate}
\end{proposition}
\begin{proof}
	(\ref{can-bij}) Let $M=X^*(T)$. Then
	\[
		\Hom(\mu_m,T)=\Hom(M,\ZZ/m)=\Hom(M,(\RR/\ZZ)[m])=\TT[m].
	\]
	This gives our desired $W$-equivariant isomorphism. Furthermore, $x\in\TT[m]$ has exact order $m$ if and only if the corresponding $\phi_x\in\Hom(\mu_m,T)$ is injective. The natural map $\Hom(\mu_m,T)\to \Hom(\mu_m,G)$ is $W$-equivariant. After fixing a choice of primitive $m$-th root of unity, \cite[Proposition 14.6]{MalleTesterman2011} proves $\Hom(\mu_m,T)/W\xrightarrow{\sim}\Conj(\mu_m,G)$.
	
	(\ref{WxNZ}) By (\ref{can-bij}), after conjugation, we may assume $\phi_x\colon\mu_m\to T$. Since $T$ is abelian, it centralizes $\phi_x$. Thus, $T\subset Z(\phi_x)$. Since $Z(\phi_x)\subset G$, this implies $T$ is a maximal torus of $Z(\phi_x)$. Next, we see $N_{Z(\phi_x)}(T)/T\subset N_G(T)/T=W$ consists of those $w\in W$ whose conjugation action on $\phi_x$ is trivial. But this is precisely the stabilizer of $x\in\TT$.
	
	(\ref{age-extension}) We see $\age_V\colon\TT\to\RR$ is well-defined and since $W$ permutes the $\beta_i$, we see $\age_V$ is constant on $W$-orbits. By definition, it is clear that $\age_V(x)=\age_V(\phi_x)$ for $x\in\TT_{\tors}$.

	(\ref{hVvsage}) We see that if $\beta_i(x)\notin\ZZ$, then $\{\beta_i(x)\}+\{-\beta_i(x)\}=1$. Otherwise, $\{\beta_i(x)\}=\{-\beta_i(x)\}=0$. Hence,
	\[
		\age_V(x) =\sum_{i=1}^s \left(\{\beta_i(x)\}+\{-\beta_i(x)\}\right)=s-\#\{i\mid \beta_i(x)=0\in\RR/\ZZ\}.
	\]
	
	(\ref{strata-cnst}) By (\ref{hVvsage}), we see $\age_\g=|\Phi^+|-h_\g$ which is integer-valued and constant on each $H_{S,I}$. We turn now to $\age_V$. Note that for any $\lambda\in X^*(T)$, the function $\{\lambda(x)\}$ is discontinuous only where it is integer-valued; thus, it is continuous on $\TT\setminus H_\lambda$ and it is also continuous on $H_\lambda$ (where it is identically $0$). Thus, $\sum_i\{\beta_i(x)\}$ is continuous on each $H_{S,I}$. Since $G\subset\SL(V)$, we see $\sum_i\beta_i(x)=0$, so 
	\[
		\sum_i\{\beta_i(x)\}=\sum_i\beta_i(x)-\sum_i\lfloor\beta_i(x)\rfloor=-\sum_i\lfloor\beta_i(x)\rfloor
	\]
	is integer-valued. Since a continuous integer-valued function is constant on connected components, the result follows.
\end{proof}


	

	
	
	
	
	
\begin{proof}[{Proof of \autoref{thm:main-McKay-connected-reductive}}]
By \autoref{thm:main-lin-red-McKay} and \autoref{prop:molien}, we see $E_{\str}(V/G)$ is given by
\[
	-t^{\dim V-\dim G+r}\lim_{q\to\infty}\sum_{m\ge1}\ \sum_{\phi\in\Conj_{\inj}(\mu_m,G)} \frac{t^{-\age_V(\phi)+\age_{\g}(\phi)}}{|W_{Z(\phi)}|}\sum_{w\in W_{Z(\phi)}}\frac{1}{\det(t-w\mid  M_\QQ)}\sum_{n\ge1}q^{nm}.
\]
Using \autoref{prop:TT-basics}(\ref{can-bij}), we may rewrite this as a sum over $\TT_{\tors}$; by \autoref{prop:TT-basics}(\ref{WxNZ}), we can then identify $W_{Z(\phi_x)}$ with the stabilizer $W_x$, and by (\ref{hVvsage}) we may write $\age_\g=|\Phi^+|-h_\g$. Thus, 
\begin{align*}
	E_{\str}(V/G) &=-t^{\dim V-|\Phi^+|}\lim_{q\to\infty}\sum_{m\ge1}\ \sum_{\substack{x\in\TT/W\\\ord(x)=m}} \frac{t^{-\age_V(x)-h_\g(x)}}{|W_x|}\sum_{w\in W_x}\frac{1}{\det(t-w)}\sum_{n\ge1}q^{nm}\\
	&=-\frac{t^{\dim V-|\Phi^+|}}{|W|}\lim_{q\to\infty}\sum_{m\ge1}\ \sum_{\substack{x\in\TT\\\ord(x)=m}} t^{-\age_V(x)-h_\g(x)}\sum_{w\in W_x}\frac{1}{\det(t-w)}\sum_{n\ge1}q^{nm}\\
	&=-\frac{t^{\dim V-|\Phi^+|}}{|W|}\lim_{q\to\infty}\sum_{n\ge1}q^n \sum_{x\in\TT[n]} t^{-\age_V(x)-h_\g(x)}\sum_{w\in W_x}\frac{1}{\det(t-w)}\\
	&=-\frac{t^{\dim V-|\Phi^+|}}{|W|}\sum_{w\in W}\frac{1}{\det(t-w)}\lim_{q\to\infty}\sum_{n\ge1}q^n \sum_{x\in\TT[n]^w} t^{-\age_V(x)-h_\g(x)}
\end{align*}
where the second equality is by the orbit-stabilizer theorem. Let $S$ be the union of $\Phi^+$ and the weights of $V$. Let $\cC$ denote the set of strata determined by $S$. By \autoref{prop:TT-basics}(\ref{strata-cnst}), $\age_V(x)+h_\g(x)$ is constant on each $C\in\cC$ with value $a(C)$. So,
\[
	\lim_{q\to\infty}\sum_{n\ge1}q^n \sum_{x\in\TT[n]^w} t^{-\age_V(x)-h_\g(x)} =
		\sum_{C\in\cC}t^{-a(C)}\lim_{q\to\infty}\sum_{n\ge1}\#(C\cap\TT^w\cap\TT[n]) q^n.
\]
We claim that 
\begin{equation}\label{eqn:HLchi}
	-\lim_{q\to\infty}\sum_{n\ge1}\#(C\cap\TT^w\cap\TT[n]) q^n=\chi_c(C\cap\TT^w).
\end{equation}
Indeed, letting $r=\rank M$, we have a quotient map $q\colon\RR^r\to\RR^r/\ZZ^r\simeq\TT$. Note that $\Delta_{C,w}=[0,1)^r\cap q^{-1}(C\cap\TT^w)$ is defined by $(w-1)x\in\ZZ^r$ as well as finitely many equations of the form $\lambda_i(x)\in\ZZ$ and finitely many equations of the form $\lambda'_j(x)\notin\ZZ$, so $\Delta_{C,w}$ is bounded and definable. Since $\#(\Delta_{C,w}\cap \frac{1}{n}\ZZ^r)=\#(C\cap\TT^w\cap\TT[n])$, applying \cite[Lemma 8.5.2]{HrushovskiLoeser2015} with $\ell=0$ shows that the lefthand side of Equation \eqref{eqn:HLchi} is equal to the $o$-minimal Euler characteristic $\chi_{\omin}(\Delta_{C,w})$.\footnote{Note that in \cite{HrushovskiLoeser2015}, they denote the $o$-minimal Euler characteristic by $\chi$.} To see this agrees with $\chi_c(C\cap\TT^w)$, we partition $[0,1)^r$ into locally closed subsets $Z_I\subset[0,1)^r$, where $I\subset\{1,\dots,r\}$ and $(x_1,\dots,x_r)\in Z_I$ if and only if $I=\{i\mid x_i=0\}$. Then $q|_{Z_I}$ is a homeomorphism onto its image, so 
\[
	\chi_c(C\cap\TT^w) = \sum_I \chi_c(Z_I\cap\Delta_{C,w}) = \sum_I \chi_{\omin}(Z_I\cap\Delta_{C,w})=\chi_{\omin}(\Delta_{C,w}),
\]
where the second equality is by \cite[Lemma 3.6]{Beke2011}.

Thus,
\begin{align*}
E_{\str}(V/G) &=\frac{t^{\dim V-|\Phi^+|}}{|W|}\sum_{w\in W}\frac{1}{\det(t-w)}\sum_{C\in\cC}t^{-a(C)}\chi_c(C\cap\TT^w)\\
&=\frac{t^{\dim V-|\Phi^+|}}{|W|}\sum_{w\in W}\frac{\int_{\TT^w}t^{-\age_V-h_\g}d\chi_c}{\det(t-w)}
\end{align*}
as desired.
\end{proof}


We end this section by noting that a modified version of the above proof also applies to (disconnected) diagonalizable groups.



\begin{proposition}\label{prop:mckay-diagonalizable}
Let $V$ be a generic representation of a diagonalizable group $G\subset\SL(V)$ with $r=\rank X^*(G)$. Then
\[
	E_{\str}(V/G)=\frac{t^{\dim V}}{(t-1)^r}\int_{\TT_G}t^{-\age_V} d\chi_c 
\]
\end{proposition}
\begin{proof}
Since $H^*(BG,\QQ)=H^*(BG^\circ,\QQ)^{\pi_0(G)}$ (see e.g., \cite[9.1.1]{DeligneHodgeIII}), we have $E(BG;t)=(t-1)^{-r}$. The rest of the proof is exactly the same as for \autoref{thm:main-McKay-connected-reductive}, except we now apply \cite[Lemma 8.5.2]{HrushovskiLoeser2015} to the $\pi_0(G)$ copies of $\TT_{G^\circ}$.
\end{proof}







\section{McKay correspondence for connected reductive groups, part 2}



In this section we prove \autoref{thm:main-McKay-connected-reductive-2}. We keep the notation introduced right before \autoref{def:specialization-mat-poly}. We let $\1_S$ denote the indicator function of a set $S$. Given a finite set $S\subset M$, we let $H_{S}:=\bigcap_{\lambda\in S}H_\lambda\subset\TT$ and $\overline{S}_w:=\pi_w(S)$.



\begin{lemma}\label{l:TwcapHS-arithmetic-Tutte}
Let $w\in W$, let $S\subset M$ be a finite subset. Then
\[
		\psi_{w,S}\colon\TT^w\cap H_{S}\xrightarrow{\simeq} \Hom( M_w/\ZZ\overline{S}_w,\RR/\ZZ),\qquad x\mapsto(\lambda\mapsto\lambda(x))
\]
is an isomorphism of topological abelian groups. Furthermore,
\[
	\chi_c(\TT^w\cap H_{S})=
	\begin{cases}
		|(M_w/\ZZ\overline{S}_w)_{\tors}|,& \dim \QQ\overline{S}_w=r_w\\
		0,&\mathrm{else}
	\end{cases}
\]
\end{lemma}
\begin{proof}
The $(w,S)=(1,\varnothing)$ case holds by definition. Then note that $x\in\TT=\Hom(M,\RR/\ZZ)$ is stabilized by $w$ if and only if the function $x$ vanishes on $(w-1)M$. Similarly, $x\in H_S$ if and only if $x$ vanishes on $\ZZ S$. Thus, 
\[
	\TT^w\cap H_S=\Hom(M/((w-1)M+\ZZ S),\RR/\ZZ)=\Hom(M_w/\ZZ\overline{S}_w,\RR/\ZZ).
\]

%It is clear that $\psi_{1,\varnothing}$ is an isomorphism. Thus, it is a continuous bijection from a compact space to a Hausdorff space, and so is a homeomorphism.
%
%We claim we have a commutative diagram
%\[
%\xymatrix{
%	\TT^w\cap H_S\ar[r]^-{\psi_{w,S}}\ar@{^(->}[d] & \Hom( M_w/\ZZ\overline{S}_w,\RR/\ZZ)\ar@{^(->}[d]\\
%	\TT\ar[r]^-{\psi_{1,\varnothing}} & \Hom( M,\RR/\ZZ)
%}
%\]
%Since
%\[
%	\Hom( M_w/\ZZ\overline{S}_w,\RR/\ZZ)=\Hom( M/((w-1) M+\ZZ S),\RR/\ZZ),
%\]
%commutativity of our diagram follows from the fact that for all $x\in\TT^w$, $\mu\in M$, and $a_\lambda\in\ZZ$, we have
%\[
%	((w-1)\mu+\sum_{\lambda\in S}a_\lambda\lambda)(x)=(w\mu)(x)-\mu(x)+\sum_{\lambda\in S}a_\lambda\lambda(x)=0;
%\]
%note $\lambda(x)=0$ since $x\in H_S\subset H_\lambda$, and $(w\mu)(x)-\mu(x)=\mu(w^{-1}x)-\mu(x)=0$ since $wx=x$.
%
%Since the diagram commutes, we see $\psi_{w,S}$ is injective. For surjectivity, given $\varphi\in\Hom( M_w/\ZZ\overline{S}_w,\RR/\ZZ)$, we know there exists $x\in\TT$ for which $\psi_{1,\varnothing}(x)=\varphi$ and we must show $x\in\TT^w$. By definition of $\varphi$, we see $\lambda(x)=0$ for all $\lambda\in S$, and $0=((w-1)\mu)(x)=\mu(w^{-1}x)-\mu(x)=\mu(w^{-1}x-x)$ for all $\mu$. Thus, $x\in H_S$ and $w^{-1}x=x$. This finishes the proof that $\psi_{w,S}$ is an isomorphism. We see $\psi_{w,S}$ is continuous since it is the restriction of $\psi_{1,\varnothing}$. Thus, it is again a continuous bijection from a compact space to a Hausdorff space, so is a homeomorphism.

Since $M_w/\ZZ\overline{S}_w\simeq\ZZ^{r_w-\rho_w(S)}\oplus(M_w/\ZZ\overline{S}_w)_{\tors}$, we see $\TT^w\cap H_S$ is $|(M_w/\ZZ\overline{S}_w)_{\tors}|$ copies of an $(\RR/\ZZ)^{r_w-\rho_w(S)}$. Our desired result then follows from the fact that $\chi_c((\RR/\ZZ)^k)$ vanishes for $k>0$ and is equal to $1$ if $k=0$.
\end{proof}



\begin{proof}[{Proof of \autoref{thm:main-McKay-connected-reductive-2}}]
For all $w\in W$ and $R\subset\Phi^+$, we have $T_{w,R}\subset G\subset\SL(V)$. Since $T_{w,R}\subset G$ is closed and $V$ is generic for $G$, it is also generic for $T_{w,R}$. Indeed, by definition, there exists a $G$-invariant open subset $U\subset V$ such that $\codim(V\setminus U)\geq2$, and every $v\in U$ has trivial $G$-stabilizer and closed $G$-orbit. Since $T_{w,R}\subset G$ is closed, we see the same properties hold for the $T_{w,R}$-action on $U$.

Note that
\[
	t^{h_\g}=t^{\sum_{\alpha\in\Phi^+} \1_{H_\alpha}} =\prod_{\alpha\in\Phi^+} t^{\1_{H_\alpha}}=\prod_{\alpha\in\Phi^+} (1+(t-1)\1_{H_\alpha})=\sum_{R\subset\Phi^+}(t-1)^{|R|}\1_{H_R}.
\]
Thus, we find
\begin{align*}
	\int_{\TT^w}t^{-\age_V-h_\g}d\chi_c &=\int_{\TT^w}t^{-\age_V}\sum_{R\subset\Phi^+}(t^{-1}-1)^{|R|}\1_{H_{R}}d\chi_c\\
	&=\sum_{R\subset\Phi^+}(t^{-1}-1)^{|R|}\int_{\TT^w\cap H_R}t^{-\age_V}d\chi_c.
\end{align*}
By \autoref{l:TwcapHS-arithmetic-Tutte}, $\TT^w\cap H_R=\TT_{T_{w,R}}$, where $T_{w,R}\subset T\subset G$ is the diagonalizable subgroup which has dimension $r_w-\rho_w(R)$. Thus, \autoref{prop:mckay-diagonalizable} shows
\begin{align*}
	\frac{\int_{\TT^w}t^{-\age_V-h_\g}d\chi_c}{\det(t-w\mid M_\QQ)} &= \frac{1}{t^{\dim V}\det(t-w\mid M_\QQ)}\sum_{R\subset\Phi^+}(t^{-1}-1)^{|R|}(t-1)^{r_w-\rho_w(R)}E_{\str}(V/T_{w,R})\\
	&= \frac{1}{t^{\dim V}\det(t-w\mid M_\QQ)}\sum_{R\subset\Phi^+}(-t^{-1})^{|R|}(t-1)^{|R|+r_w-\rho_w(R)}E_{\str}(V/T_{w,R})\\
	&= \frac{1}{t^{\dim V}\Delta_w(t)}\sum_{R\subset\Phi^+}(-t^{-1})^{|R|}(t-1)^{|R|-\rho_w(R)}E_{\str}(V/T_{w,R}).
\end{align*}

The assertion about $E_{\str}(V/G)$ is then immediate from \autoref{thm:main-McKay-connected-reductive}. Lastly, the computation of the stringy Euler characteristic follows from the fact that $\chi_{\str}(V/G)=\lim_{t\to1}E_{\str}(V/G;t)$.
\end{proof}











\section{McKay correspondence for symmetric representations of reductive groups}


In this section we prove \autoref{prop:main-intTw-arithmetic-Tutte} and \autoref{thm:main-McKay-connected-reductive-symmetric}. We keep the notation introduced right before \autoref{def:specialization-mat-poly}.



\begin{proof}[{Proof of \autoref{prop:main-intTw-arithmetic-Tutte}}]
Note that
\[
	t^{h_V}=t^{\sum_i \1_{H_{\beta_i}}} =\prod_i t^{\1_{H_{\beta_i}}}=\prod_i (1+(t-1)\1_{H_{\beta_i}})=\sum_{B\subset\cW^+}(t-1)^{|B|}\1_{H_B}.
\]
Writing $t^{h_\g}$ in a similar manner, we see
\begin{align*}
	\int_{\TT^w}t^{h_V-h_\g}d\chi_c &=\int_{\TT^w}\sum_{\substack{B\subset\cW^+\\ R\subset\Phi^+}}(t-1)^{|B|}(t^{-1}-1)^{|R|}\1_{H_{B\cup R}}d\chi_c\\
	&=\sum_{\substack{B\subset\cW^+\\ R\subset\Phi^+}}(t-1)^{|B|}(t^{-1}-1)^{|R|}\chi_c(\TT^w\cap H_{B\cup R})\\
	&=\sum_{\substack{B\subset\cW^+,\, R\subset\Phi^+\\ \dim\QQ(\overline{B}_w\cup \overline{R}_w)=r_w}}|( M_w/\ZZ(\overline{B}_w\cup\overline{R}_w))_{\tors}|\,(t-1)^{|B|}(t^{-1}-1)^{|R|}\\
	&=\sum_{\substack{B\subset\cW^+,\, R\subset\Phi^+\\ \dim\QQ(\overline{B}_w\cup \overline{R}_w)=r_w}}m_w(B,R)\,(t-1)^{|B|}(t^{-1}-1)^{|R|}
\end{align*}
where the third equality is by \autoref{l:TwcapHS-arithmetic-Tutte}. We see that the final expression is precisely $\cM_w(0,t-1,t^{-1}-1)$. Then using \autoref{prop:TT-basics}(\ref{hVvsage}), we see
\[
	\int_{\TT^w}t^{-\age_V-h_\g}d\chi_c=t^{-|\cW^+|}\int_{\TT^w}t^{h_V-h_\g}d\chi_c,
\]
which proves the first claim in \autoref{prop:main-intTw-arithmetic-Tutte}.

Next, notice that in $\ZZ[t^\pm]$ we may write $t^{-1}-1=-t^{-1}(t-1)$, and so
\begin{equation}\label{eqn:Mw-alternative-expression}
	\sum_{\substack{B\subset\cW^+,\, R\subset\Phi^+\\ \dim\QQ(\overline{B}_w\cup \overline{R}_w)=r_w}}m_w(B,R)\,(t-1)^{|B|}(t^{-1}-1)^{|R|}=\sum_{\substack{B\subset\cW^+,\, R\subset\Phi^+\\ \dim\QQ(\overline{B}_w\cup \overline{R}_w)=r_w}}(-t^{-1})^{|R|}m_w(B,R)\,(t-1)^{|B|+|R|}.
\end{equation}
In order for $\overline{B}_w\cup \overline{R}_w$ to span the $r_w$-dimensional vector space $( M_w)_\QQ$, we must have $|B|+|R|\geq r_w$, which proves that $(t-1)^{r_w}$ divides $\cM_w(0,t-1,t^{-1}-1)$.
\end{proof}

\begin{proof}[{Proof of \autoref{thm:main-McKay-connected-reductive-symmetric}}]
The assertion about $E_{\str}$ is immediate from \autoref{thm:main-McKay-connected-reductive} and \autoref{prop:main-intTw-arithmetic-Tutte}. The computation of $\chi_{\str}(V/G)$ follows from the fact that it is equal to $\lim_{t\to1}E_{\str}(V/G;t)$. Lastly, the expression for $\widetilde{\cM}_w(1)$ follows from \eqref{eqn:Mw-alternative-expression}.
\end{proof}









	
	
	
	
	
	
	
	
	
	
	
	
	
	
	
	
	
	
	
	
	
	
\section{Purity of crepant resolutions of reductive quotients}


We begin with the following result of independent interest.

\begin{proposition}\label{prop:crepant-lift-action}
Let $H$ be a connected smooth affine algebraic group over an algebraically closed field $k$ of characteristic $0$. Let $Y$ be a finitely presented affine scheme over $k$, and let $f\colon X\to Y$ be a projective crepant resolution. If $H$ acts on $Y$, then $X$ admits a $H$-action making $f$ equivariant.
\end{proposition}
\begin{proof}
By \cite[Theorem 1.1]{BierstoneMilmanTemkin2011}, there is a strong resolution $f'\colon X'\to Y$ which is functorial with respect to smooth morphisms. Thus, the $H$-action on $Y$ lifts to a $H$-action on $X'$. Since $f$ is crepant, Proposition 14.3 and Lemma 14.6 of \cite{SatrianoUsatine9} finish the proof.
\end{proof}



\begin{proof}[{Proof of \autoref{thm:purity-Hodge-Tate}}]
Note that the $G$-action on $V$ commutes with the action of the scalars $\bG_m\subset\GL(V)$ on $V$. Thus, the $\bG_m$-action descends to $Y=V/G$. Let $q\colon V\to Y$ be the quotient map. Since $V^{\bG_m}=\{0\}$ and $0$ in the $\bG_m$-orbit closure of every point of $V$, we see $Y^{\bG_m}=\{q(0)\}$. Indeed, if $y\in Y(\CC)$ is a fixed point, choosing $v\in V(\CC)$ with $y=q(v)$, we note $y=ty=tq(v)=q(tv)$ for all $t\in\bG_m(\CC)$; letting $t\to0$, we see $y=q(0)$.

With this in place, the fact that $H^i_c(X,\QQ)$ is pure of weight $i$ follows from same argument in \cite[Theorem 14.1]{SatrianoUsatine9} with two modifications:~(i) we use \autoref{prop:crepant-lift-action} in place of \cite[Proposition 14.5]{SatrianoUsatine9}, and (ii) we replace all instances of $\ell$-adic cohomology with compactly supported cohomology equipped with their mixed Hodge structures. By Poincar\'e duality, $H^i(X,\QQ)$ is pure of weight $i$.

Lastly, since $\pi$ is crepant, we have $E(X;u,v)=E_{\str}(Y;u,v)$. By \cite[Corollary 7.2]{HuangSatrianoUsatine}, 
\[
	\chi_{\Hdg}(\{BH\})=\sum_{i\in\ZZ}(-1)^i[H_c^i(BH,\QQ)]
\]
for all linearly reductive groups $H$. Thus, by \autoref{prop:molien} and \autoref{thm:main-lin-red-McKay}, $E(X;u,v)=E_{\str}(Y;u,v)$ is a function of $t=uv$. Since $H^i_c(X,\QQ)$ is pure of weight $i$, it follows it is Hodge--Tate. By Poincar\'e duality, $H^*(X,\QQ)$ is as well.
\end{proof}




	
	
	
	
	
	
	
	
	
	
	
	
	
	
	
	


\section{Counter-examples to \autoref{conj:main1}, \autoref{conj:main2}, and \autoref{conj:main3}}


\begin{proposition}\label{prop:mab=420-has-NCCR}
	With notation and hypotheses as in \autoref{thm:NCCR-counter-examples}, $V$ is a symmetric generic representation of $\GL_2$. Furthermore, if $(m,a,b)=(4,2,0)$, then $V/\GL_2$ has an NCCR.
\end{proposition}
\begin{proof}
	The first paragraph of the proof of \cite[Theorem 5.2.1]{SpenkoVanDenBergh17} shows $U^{\oplus m}$ is generic for $m\geq3$ (for them $m=h$ and $n=2$). We see $V$ is symmetric by definition.
	
	To show $V$ has an NCCR, we apply \cite[Theorem 1.6.4]{SpenkoVanDenBergh17} with $\mu=0$. Here $W=\ZZ/2$ acts on $X^*(T)=\ZZ^2$ by swapping $e_1$ and $e_2$. We see $\Phi^+=\{\alpha\}$ with $\alpha=e_1-e_2$. We must check condition (1.5) of (loc.~cit). We choose $\epsilon=(1,1)$ which is $W$-invariant. An explicit computation shows that $\overline{\Sigma}$ is the convex hull of $\pm(7,7)$, $\pm(7,-5)$, and $\pm(5,-7)$. Thus, $\overline{\Sigma}_{\pm\epsilon}\setminus\Sigma$ consists of the open edge connecting $(7,-5)$ and $(5,-7)$, as well as the open edge connecting $(-7,5)$ and $(-5,7)$. These are the segments $\{(6,-6) + t(1,1) \mid t\in(-1,1)\}$ and $\{(-6,6) + t(1,1) \mid t\in(-1,1)\}$. Thus, 
\begin{align*}
	-\overline{\rho}+\frac{1}{2}&(\overline{\Sigma}_{\pm\epsilon}\setminus\Sigma) \\
	 &= \{(2.5,-2.5) + t(1,1) \mid t\in(-0.5,0.5)\} \cup \{(-3.5,3.5) + t(1,1) \mid t\in(-0.5,0.5)\},
\end{align*}
which contains no lattice points. Thus, \cite[Theorem 1.6.4]{SpenkoVanDenBergh17} shows that $V/\GL_2$ has an NCCR.
\end{proof}




\begin{proof}[{Proof of \autoref{thm:NCCR-counter-examples}}]
In light of \autoref{prop:mab=420-has-NCCR}, we need only prove the claim about $\chi_{\str}(V/\GL_2)$. To begin, choose a linear functional $\ell\colon X^*(T)=\ZZ^2\to\ZZ$ to be $\ell(x,y)=x+y$; using $\ell$, the multi-set of positive weights $\cW^+$ of $V$ is given by $e_1=(1,0)$ with multiplicity $m$, $e_2=(0,1)$ with multiplicity $m$,
\[
	\beta_i:=(a-i,b+i),\qquad 0\leq i\leq a-b
\]
each with multiplicity $1$. We have $\Phi^+=\{\alpha\}$ with $\alpha=e_1-e_2$. Let $s$ be the non-trivial generator of $W\simeq\ZZ/2$.

By \autoref{thm:main-McKay-connected-reductive-symmetric}, we have
\[
	\chi_{\str}(V/\GL_2)=\frac{1}{2}\left(\widetilde{\cM}_1(1)+\frac{\widetilde\cM_s(1)}{2}\right).
\]
For $w\in W$, we compute $\widetilde{\cM}_w(1)$ using \autoref{thm:main-McKay-connected-reductive-symmetric}.

We begin with $\widetilde{\cM}_s(1)$. Note that $\pi_s\colon\ZZ^2\simeq M\to M_s=M/\ZZ\alpha\simeq\ZZ$ is given by $(x,y)\mapsto x+y$. Since $r_s=1$, we see
\[
	\widetilde{\cM}_s(1) =\sum_{\lambda\in\cW^+}m_s(\{\lambda\},\varnothing)=2m+(a+b)(a-b+1),
\]
where the $2m$ term is from $\lambda=e_i$, and the $(a+b)(a-b+1)$ term is from the remaining positive weights.


We next calculate $\widetilde{\cM}_1(1)$. Since $M_1=M=\ZZ^2$ and $r_1=2$, \autoref{thm:main-McKay-connected-reductive-symmetric} tells us
\begin{align*}
\widetilde{\cM}_1(1) &= \sum_{\substack{B\subset\cW^+\\ |B|=\rho_1(B)=2}}m_1(B,\varnothing)-\sum_{\substack{\beta\in\cW^+\\\rho_1(\{\beta,\alpha\})=2}}m_1(\{\beta\},\{\alpha\})\\
&=\sum_{\substack{\{\lambda,\mu\}\subset\cW^+\\ \rho_1(\{\lambda,\mu\})=2}}|\det(\lambda,\mu)|-\sum_{\lambda\in\cW^+}|\det(\lambda,\alpha)|.
\end{align*}
Since 
\[
|\det(e_1,\alpha)|=|\det(e_2,\alpha)|=1,\qquad |\det(\beta_i,\alpha)|=a+b,
\]
we see
\[
\sum_{\lambda\in\cW^+}|\det(\lambda,\alpha)|=2m+(a+b)(a-b+1).
\]
Next, since
\[
	|\det(e_1,\beta_i)|+|\det(e_2,\beta_i)|=(b+i)+(a-i)=a+b
\]
and
\[
\sum_{0\leq i<j\leq n}|\det(\beta_i,\beta_j)|=(a+b)\sum_{0\leq i<j\leq n}(j-i)=(a+b)\binom{a-b+2}{3},
\]
we see
\[
	\sum_{\{\lambda,\mu\}\subset\cW^+}|\det(\lambda,\mu)|=m^2+m(a+b)(a-b+1)+(a+b)\binom{a-b+2}{3},
\]
and hence
\[
	\widetilde{\cM}_1(1)=m^2-2m+(a+b)(a-b+1)(m-1)+(a+b)\binom{a-b+2}{3}.
\]

Combining these expressions, we find
\[
	\chi_{\str}(V/\GL_2)=\frac{m(m-1)}{2}+\frac{m(a+b)(a-b+1)}{2}+\frac{(a+b)(a-b-1)(a-b+1)(a-b+3)}{12}.
\]
The first term is an integer, and the second term is as well since $a+b$ and $a-b+1$ have different parities. 

If $a-b=2k+1$ is odd, then
\[
	\frac{(a-b-1)(a-b+1)(a-b+3)}{12}=\frac{2k(k+1)(k+2)}{3}=4{k+2\choose 3}\in\ZZ.
\]

If $a-b=2k$ is even, then $(a-b-1)(a-b+1)(a-b+3)\equiv(-1)(1)(3)\equiv3\pmod{6}$, so $\frac{1}{3}(a-b-1)(a-b+1)(a-b+3)$ is an odd integer. Thus, $\chi_{\str}(V/\GL_2)\not\in\ZZ$ if and only if $\frac{1}{4}(a+b)\not\in\ZZ$. Since $a+b$ is even, this is equivalent to $a+b\equiv2\pmod{4}$.
\end{proof}














\bibliographystyle{alpha}
\bibliography{McKay}




\end{document}



